% ================================================================
\subsection{Gravitational Universality: Every Coupling is Attractive}
\label{sec:gravitational-universality}
% ================================================================

The Gravitational Universality Theorem
(\lean{GravitationalUniversality.lean}, 3~theorems, zero sorry,
zero custom axioms, proved July~16, 2026) establishes the strongest
possible positivity result for the Gram correlator: not merely that the
matrix $G$ is positive definite (a collective property), but that
\emph{every individual entry} is strictly positive.

\begin{theorem}[Gravitational Universality (Proved)]
\label{thm:gravitational-universality}
For all $j, k \geq 1$:
\[
  G(j,k) = \int_0^1 \fract{1/(jx)}\,\fract{1/(kx)}\,dx > 0.
\]
\end{theorem}

In the physical dictionary, this says:
\begin{quote}
\emph{Every mode couples positively to every other mode.
The prime number quantum field is universally attractive---there
are no repulsive interactions.}
\end{quote}

This is the number-theoretic incarnation of the universality of
gravitation: just as every mass attracts every other mass, every
integer-indexed basis function $h_k(x) = \fract{1/(kx)}$ has
strictly positive inner product with every other $h_j$.

\subsubsection{Proof Architecture}

The proof decomposes into three cases:

\begin{enumerate}
  \item \textbf{Diagonal ($j = k$):} $G(k,k) = k^{-2}(\ln 2\pi - \gamma - 1) > 0$.
    Already proved in \lean{vasyuninGramEntry\_diag\_pos}
    (Structural.lean).

  \item \textbf{Off-diagonal, $j \mid k$:} The Vasyunin integral identity
    (\lean{vasyunin\_eq\_integral\_proved}, zero axioms) rewrites
    $G(j,k)$ as the integral $\int_0^1 \fract{1/(jx)} \cdot
    \fract{1/(kx)}\,dx$. On the interval $(1/(j+k),\; 1/k)$:
    \begin{itemize}
      \item $1/(kx) \in (1,2)$, so $\fract{1/(kx)} = 1/(kx) - 1 > 0$.
      \item Writing $k = jq$, we have $1/(jx) \in (q, q+1)$
        with no integer inside the open interval, so
        $\fract{1/(jx)} > 0$.
    \end{itemize}
    The integrand is pointwise positive on $(1/(j+k), 1/k)$.
    Since the integrand is nonneg everywhere and strictly positive
    on a subinterval, $\int_0^1 > 0$ by
    \lean{intervalIntegral\_pos\_of\_pos\_on}.

  \item \textbf{Off-diagonal, $j \nmid k$:} The integer
    $m = \lceil k/j \rceil$ creates a single zero of
    $\fract{1/(jx)}$ at $x_0 = 1/(jm)$. We avoid it by restricting
    to $( 1/(jm),\; 1/k )$, where $1/(jx) \in (k/j,\; m)$ contains
    no integer. The same $\fract{1/(kx)} > 0$ bound applies,
    giving pointwise positivity.
\end{enumerate}

The case $j > k$ follows by the symmetry $G(j,k) = G(k,j)$
(commutativity of multiplication in the integrand).

\begin{correspondence}[Gravity vs.\ Electromagnetism]
In the Standard Model of particle physics, the electromagnetic
force can be attractive or repulsive (opposite charges attract,
like charges repel), while gravity is \emph{universally attractive}.
The Gram correlator exhibits gravitational behavior: $G(j,k) > 0$
for \emph{all} pairs, with no sign changes. This is consistent with
the spectral observation that $G$ follows GOE (Gaussian Orthogonal
Ensemble) statistics rather than GUE---a real symmetric positive
operator, not a complex Hermitian one with possible sign oscillations.

In the lattice QFT language of \S\ref{sec:vasyunin}, gravitational
universality says that the partition function of any pair of
lattice sites is strictly positive: there are no frustration effects,
no antiferromagnetic couplings, no sign problem. The prime number
quantum field is a \emph{purely ferromagnetic} system.
\end{correspondence}

\begin{remark}[Subinterval Analysis as Microlocal Technique]
The proof technique---constructing specific subintervals where both
fractional parts are simultaneously bounded away from zero---is the
number-theoretic analogue of \emph{microlocal analysis} in PDE theory.
Rather than proving global positivity of the integrand (which is false:
$\fract{1/(jx)} = 0$ at $x = 1/(jn)$ for every integer $n$), the proof
identifies a ``microsupport''---a specific region of phase space where
the wavefront set is controlled---and exploits the strict positivity
there. The case split on $j \mid k$ vs.\ $j \nmid k$ is a partition
of unity adapted to the arithmetic geometry of the fractional parts.
\end{remark}
